% 第九卷 6.5.4.1 节：EAGLE 草稿器的结构。
% 目标模型冻结；草稿器 = 拼接 + 线性层（2d→d）+ 一个解码器层；嵌入层与输出头取自目标模型。
\begin{tikzpicture}[
  x=1cm,y=1cm,
  every node/.style={font=\figurefont\footnotesize,text=FigureInk},
  frozen/.style={rectangle,rounded corners=1.8pt,draw=FigureStroke,line width=0.6pt,fill=FigurePaper,minimum height=8mm,align=center,inner sep=2pt},
  draft/.style={rectangle,rounded corners=1.8pt,draw=FigureYellow,line width=0.6pt,fill=FigureYellowFill,minimum height=8mm,align=center,inner sep=2pt},
  data/.style={rectangle,rounded corners=1.8pt,draw=FigureBlue,line width=0.6pt,fill=FigureBlueFill,minimum height=6mm,align=center,inner sep=2pt},
  arr/.style={-{Stealth[length=1.8mm]},FigureStroke,line width=0.6pt},
  lab/.style={font=\figurefont\scriptsize,text=FigureMuted}
]
  \node[frozen,minimum width=20mm] (T) at (0,0.75) {目标模型\\（冻结）};
  \node[data] (F) at (2.75,0.75) {$\bm f_1,\ldots,\bm f_t$};
  \node[data] (X) at (0,-0.75) {已选词元\\$x_2,\ldots,x_{t+1}$};
  \node[frozen] (E) at (2.75,-0.75) {嵌入层\\（共享）};
  \node[draft] (C) at (5.3,0) {拼接\\线性层$2d\to d$};
  \node[draft] (Dc) at (7.55,0) {解码器层};
  \node[data] (Fh) at (9.5,0) {$\hat{\bm f}_{t+1}$};
  \node[frozen] (O) at (11.45,0) {输出头\\（共享）};
  \node[data] (P) at (13.6,0) {$\hat p_{t+2}$\\抽取$\hat x_{t+2}$};
  \draw[arr] (T) -- (F);
  \draw[arr] (X) -- (E);
  \draw[arr] (F.east) -- ++(0.35,0) |- ([yshift=1.5mm]C.west);
  \draw[arr] (E.east) -- ++(0.35,0) |- ([yshift=-1.5mm]C.west);
  \draw[arr] (C) -- (Dc);
  \draw[arr] (Dc) -- (Fh);
  \draw[arr] (Fh) -- (O);
  \draw[arr] (O) -- (P);
  \draw[arr,dashed,FigureMuted] (P.south) -- ++(0,-1.05) -| (C.south);
  \node[lab,anchor=north] at (9.5,-1.6) {$(\hat{\bm f}_{t+1},\hat x_{t+2})$接到输入末尾，生成下一个草稿};
  \node[lab,anchor=south] at (6.3,0.62) {草稿器（可训练）};
  \draw[FigureYellow,dashed,line width=0.5pt,rounded corners=2pt] (4.1,-0.6) rectangle (8.45,1.1);
\end{tikzpicture}
